\chapter{Вакуумный отклик на неизбежный дефицит замыкания}

Предыдущая глава доказала, что внутри класса пятнадцать совершенное
замыкание невозможно. Теперь детерминант рассматривается в правильном
порядке причинности: не как источник дефекта, а как отклик вакуума на уже
вынужденный компактный проектор.

В отличие от общих абсолютных детерминантов, здесь сравниваются два
оператора, различающиеся на конечномерном подпространстве. Поэтому
относительный ответ конечен и не требует подбора локального контрчлена.

\section{Замкнутый и дефектный операторы}

На одной ориентированной ветви положим
\begin{equation}
 H_{\rm cl}=I,
 \qquad
 H_{\rm def}=\widehat T_+\widehat T_+^*=I-P_{\rm def},
 \qquad
 P_{\rm def}=|e_0\rangle\langle e_0|\otimes q_0.
 \label{eq:v5-closed-defect-hamiltonians}
\end{equation}
Проектор $P_{\rm def}$ имеет ранг пятнадцать. Функциональное исчисление
даёт точное равенство
\begin{equation}
 e^{-tH_{\rm def}}-e^{-tH_{\rm cl}}
 =(1-e^{-t})P_{\rm def}.
 \label{eq:v5-finite-rank-heat-response}
\end{equation}
Следовательно, нормированный относительный тепловой след равен
\begin{equation}
 K_{\rm def}(t)
 =\frac1{105}\Tr
 \left(e^{-tH_{\rm def}}-e^{-tH_{\rm cl}}\right)
 =\frac17(1-e^{-t}).
 \label{eq:v5-defect-heat-trace}
\end{equation}

При $t\to0$ ответ стремится к нулю. Ультрафиолетовые локальные коэффициенты
двух операторов сокращаются. При $t\to\infty$ остаётся $1/7$: длинноволновый
отклик видит ровно пятнадцать нулевых каналов. Это является конечномерной
версией общего heat-kernel чтения индекса; связь суперследа, нулевых мод и
индекса обсуждается, например, в \path{arXiv:math/0512167} и
\path{arXiv:math/0304214}.

\section{Точный относительный детерминант}

Введём положительный регулятор или физическую щель $m>0$ и определим
отклик относительно гипотетически замкнутого оператора:
\begin{equation}
 \Gamma_{\rm def}(m)
 =\int_0^\infty\frac{dt}{t}e^{-m^2t}K_{\rm def}(t).
 \label{eq:v5-defect-proper-time-response}
\end{equation}
Интеграл сходится на обоих концах. По формуле Фруллани
\begin{equation}
 \Gamma_{\rm def}(m)
 =\frac17\log\frac{m^2+1}{m^2}.
 \label{eq:v5-defect-relative-determinant}
\end{equation}
Тот же результат получается непосредственно из конечномерного
фредгольмова детерминанта: на $P_{\rm def}$ собственное значение меняется
с $m^2+1$ на $m^2$, а на ортогональном дополнении всё сокращается.

Таким образом, топология действительно порождает универсальный
\emph{относительный} вакуумный отклик. Его кратность полностью фиксирована
и равна $1/7$. Это сильнее общего утверждения о допустимости фермионного
детерминанта: локальная схема здесь не входит в ответ.

\section{Почему отклик не выбирает масштаб}

Число $m$ не выводится индексом. В безмассовом пределе
\begin{equation}
 \Gamma_{\rm def}(m)\longrightarrow+\infty,
 \qquad m\to0,
\end{equation}
что лишь регистрирует точные нулевые моды. При $m\to\infty$ дефектный
отклик исчезает.

Если ненулевой уровень имеет величину $a>0$, то
\begin{equation}
 \Gamma_{\rm def}(m,a)
 =\frac17\log\left(1+\frac{a}{m^2}\right).
 \label{eq:v5-defect-response-general-gap}
\end{equation}
Для внутреннего масштабирования $a(R)=R^{-2}$ функция строго убывает по
$R$ и не имеет конечной стационарной точки. Она предпочитает размазывание
дефекта к $R\to\infty$, а не локализованный радиус.

\begin{theorem}[Универсальный отклик без выбора размера]
Компактный дефект класса пятнадцать имеет конечный scheme-independent
относительный детерминант с точной кратностью $1/7$. Но ответ зависит от
отношения ненулевого уровня к щели $m$ и монотонен при масштабировании
$a=R^{-2}$. Поэтому индекс порождает вакуумный отклик, но не определяет
массу, конечный радиус или абсолютную энергию.
\end{theorem}

\section{Обменная Real-пара и два разных чтения}

Для сопряжённой ветви ориентированный индекс меняет знак. Поэтому
индексный суперслед двух ветвей сокращается. Но положительные дефектные
проекторы присутствуют в обеих половинах. Сравнение модуля полного
квадратичного оператора с замкнутым эталоном даёт
\begin{equation}
 \frac{30}{210}
 \log\left(1+\frac{a}{m^2}\right)
 =\frac17\log\left(1+\frac{a}{m^2}\right).
\end{equation}

Следовательно, необходимо различать:
\begin{enumerate}
 \item ориентированный суперслед, который хранит знак индекса и
 сокращается в Real-паре;
 \item положительный модуль детерминанта, который считает оба дефектных
 подпространства и сохраняет вес $1/7$.
\end{enumerate}
Выбор между фазой, суперследом и положительным модулем является вопросом
физической меры. Старые гейты Pfaffian и determinant line уже показали,
что этот выбор нельзя делать вручную. Геометрия determinant line и её
связь с эта-инвариантом описаны в \path{arXiv:dg-ga/9505002}; для
представителей $K^1$ и связи с функционалом Весса--Зумино см.
\path{arXiv:1205.0562}.

\begin{nogocriterion}[Запрет превращения отклика в абсолютную энергию]
Нельзя объявлять \eqref{eq:v5-defect-relative-determinant} массой частицы,
полагая $m=1$ или отбрасывая логарифм. Значение $m$ и выбор физического
чтения determinant line должны следовать из родительской меры до
сравнения с наблюдаемыми.
\end{nogocriterion}

\section{Новая причинная цепочка}

Получена строгая последовательность
\begin{equation}
 \text{ненулевой KO-класс}
 \Longrightarrow
 \text{неизбежный дефект}
 \Longrightarrow
 \text{конечный относительный отклик}.
\end{equation}
Обратная последовательность, в которой детерминант сначала выбирает
дефект, больше не требуется. Открыты два последних звена: почему вакуум
выбирает ненулевой класс и какая аномальная фаза Real-пары задаёт
физическую ориентацию меры.

\begin{protocol}
\begin{enumerate}
 \item разность тепловых ядер: \textbf{конечного ранга};
 \item локальные ультрафиолетовые члены: \textbf{сокращаются};
 \item инфракрасный предел: \textbf{$1/7$};
 \item относительный детерминант: \textbf{вычислен точно};
 \item зависимость от щели $m$: \textbf{остаётся};
 \item конечный стационарный радиус: \textbf{нет};
 \item положительный отклик Real-пары: \textbf{сохраняет $1/7$};
 \item ориентированный суперслед Real-пары: \textbf{сокращается};
 \item абсолютная энергия: \textbf{не выведена};
 \item физическое замыкание: \textbf{не достигнуто};
 \item следующий гейт: \textbf{эта-фаза и ориентация Real-пары}.
\end{enumerate}
\end{protocol}

Машинный сертификат сохранён в
\path{s2t_v5_closure_deficit_induced_vacuum_response_gate_results.json}.